\chapter{Parte I: transferencia de Stockformer al S\&P 500}\label{cap:parte1}

Esta parte responde a la primera pregunta de investigación, esto es, si Stockformer, un modelo de
selección de acciones basado en wavelets, atención y grafos de relaciones, calibrado sobre
acciones chinas de clase A, conserva su capacidad predictiva sobre el universo del S\&P
500 descrito en la metodología. En vez de comparar Stockformer contra un único rival, se
le sitúa dentro de la escalera de complejidad de la sección~\ref{sec:met-ladder}, que
ordena los modelos desde controles sin parámetros entrenables hasta la red de en torno a
un millón de parámetros.

\section{Resultados de la escalera de complejidad}\label{sec:p1-ladder}

La escalera evalúa, sobre un mismo panel y un mismo marco de evaluación, una familia de
modelos ordenada por número de parámetros y capacidad. El procedimiento es el mismo para
todos. Las predicciones se convierten en una ordenación transversal diaria, se calcula el
coeficiente de información de rangos (IC) sobre la ventana de prueba y, en paralelo, se
construye una cartera de posiciones largas y cortas para medir el Sharpe asociado. Un IC
positivo y estable es condición necesaria para que un modelo aporte señal. La
Tabla~\ref{tab:ladder} recoge los resultados; como la ventana efectiva de Stockformer (236
días) difiere de la del resto (250 días), las comparaciones directas entre modelos se
hacen en la Sección~\ref{sec:p1-signif} sobre los días comunes.

\begin{table}[htbp]
\caption{Complexity ladder: IC, parámetros y Sharpe por modelo (split único, 250 días).}
\label{tab:ladder}
\begin{tabular}{llrrrrrr}
\toprule
model & family & n\_params & ic\_mean & icir & tstat & pvalue & sharpe \\
\midrule
zero & sanity & 0 &  &  &  &  &  \\
momentum & sanity & 0 & -0.0128 & -0.0670 & -1.0589 & 0.2907 & -1.4289 \\
reversal & sanity & 0 & 0.0126 & 0.0673 & 1.0635 & 0.2886 & -1.4258 \\
ridge & linear & 376 & 0.0173 & 0.1064 & 1.6826 & 0.0937 & -2.8542 \\
lasso & linear & 5 & 0.0238 & 0.0951 & 1.5035 & 0.1340 & 0.9099 \\
elasticnet & linear & 19 & 0.0157 & 0.0687 & 1.0864 & 0.2784 & -0.0607 \\
lightgbm & trees & 31 & 0.0136 & 0.1110 & 1.7547 & 0.0805 & 0.2062 \\
xgboost & trees & 9600 & 0.0141 & 0.0926 & 1.4639 & 0.1445 & -1.8505 \\
\bottomrule
\end{tabular}
\end{table}


La escalera se recorre de abajo arriba. En el peldaño cero están los controles sin
parámetros entrenables: la cartera nula, el momentum y la reversión. Sirven para calibrar
qué valor de IC es indistinguible del ruido en este panel. El momentum obtiene un IC de
$-0{,}0128$ ($t=-1{,}06$, $p=0{,}29$) y la reversión un IC de $+0{,}0126$ ($t=1{,}06$,
$p=0{,}29$); ambos son simétricos y no significativos. Esa simetría fija una banda de
en torno a $\pm 0{,}013$ como magnitud de lo que el azar produce sobre 250 días.

En el primer peldaño (modelos lineales) está el mejor modelo de la escalera. El Lasso, con
5 parámetros efectivos tras la regularización, alcanza un IC de $+0{,}0238$ y el mayor
Sharpe de la tabla en esta partición, de $+0{,}91$ (el único otro positivo es el de
LightGBM, $+0{,}21$), con un retorno anualizado del $27{,}4\%$. Esta última cifra sale de
una única partición de 250 días y es un diagnóstico de la señal, no un resultado de
cartera; el resultado de cartera defendible es el del protocolo de avance temporal del
Capítulo~\ref{cap:parte2}. El Ridge, con 376 parámetros, obtiene un IC de $+0{,}0173$ pero
un Sharpe de $-2{,}85$, ya que predice direccionalmente algo mejor que el ruido aunque su cartera
no es rentable, lo que muestra que IC y Sharpe no son intercambiables. La red elástica
queda en un IC de $+0{,}0157$, en línea con el resto de lineales.

En el segundo peldaño (árboles) está LightGBM, con un IC de $+0{,}0136$ y un Sharpe de
$+0{,}21$. Su capacidad efectiva en la escalera es de solo 31 hojas, porque, aunque se le
permitían hasta 300 iteraciones, la parada temprana sobre el conjunto de validación
detuvo la potenciación en la primera, de modo que el modelo se reduce a un único árbol de
31 hojas, pues el algoritmo no encuentra estructura que añadir más allá del primer árbol.
XGBoost (9.600 hojas nominales) obtiene un IC de $+0{,}0141$ y un Sharpe de $-1{,}85$. Los
árboles no mejoran a los lineales, pues su IC es parecido o inferior al del Lasso pese a tener
mucha más capacidad, y el salto de tres órdenes de magnitud en parámetros entre LightGBM y
XGBoost no da ganancia de IC.

En el peldaño superior está Stockformer, con del orden de un millón de parámetros. Su IC
es de $-0{,}003$, negativo y estadísticamente indistinguible de cero. En esta partición,
el modelo más complejo de la escalera no supera a un Lasso de cinco parámetros y queda por
debajo de la banda de ruido que fijan los controles.

\begin{figure}[htbp]
  \centering
  \includegraphics[width=0.85\linewidth]{complexity_vs_ic}
  \caption{IC medio sobre la ventana de prueba frente al número de parámetros de cada
  modelo de la escalera de complejidad, con el eje horizontal en escala logarítmica. Los
  controles sin parámetros entrenables (momentum y reversión) se sitúan por convención en
  $x=1$, al no ser representable el cero en escala logarítmica. El IC se calcula como
  media de la serie diaria de correlaciones de Spearman descrita en la
  sección~\ref{sec:met-harness}; ventana de prueba de 250 días (236 para Stockformer).}
  \label{fig:complexity-vs-ic}
\end{figure}

La Figura~\ref{fig:complexity-vs-ic} resume el resultado. Si el poder predictivo creciera
con la capacidad del modelo, cabría esperar una nube de puntos ascendente al moverse hacia
la derecha en el eje de parámetros. Lo que se observa es una relación plana en la zona de
los modelos simples y un descenso hacia Stockformer. En este panel, el IC no aumenta con
la complejidad: el mejor modelo es el más simple y el peor es el más expresivo. Este
patrón coincide, en la selección transversal de acciones, con el que
\textcite{gu2020empirical} documentan en la valoración de activos con aprendizaje
automático.

\section{Significancia estadística}\label{sec:p1-signif}

Conviene separar dos preguntas con soporte estadístico distinto. La primera es si
Stockformer muestra señal predictiva en este diseño. Su IC es de $-0{,}003$,
indistinguible de cero, frente a los IC del orden de $0{,}06$ que las arquitecturas
profundas reportan en su mercado de origen. En esta implementación, por tanto, el poder
predictivo de origen no se reproduce, y eso se sostiene contra el cero e incluso
concediéndole la ventaja del grafo (sección~\ref{sec:met-leakage}); qué parte de esa
ausencia corresponde al modelo y qué parte a la reimplementación es una cuestión que el
diseño no puede separar (sección~\ref{sec:met-amenazas}). El resultado de la escalera (que
el IC no crece con la complejidad y que el modelo más complejo es el peor) es una
afirmación sobre la forma de la relación observada en los datos. La segunda pregunta, si
los modelos simples baten a Stockformer de forma significativa, es más exigente, y es la
que esta sección mide con cuidado para no afirmar más de lo que la evidencia sostiene. La
Tabla~\ref{tab:significancia} recoge, para cada modelo simple, la diferencia de IC frente
a Stockformer, el estadístico $t$ del contraste pareado por día y su $p$-valor, sobre
$n=236$ días emparejados.

\begin{table}[htbp]
\caption{Tests de significancia del IC frente a Stockformer (diferencia de medias, t, p).}
\label{tab:significancia}
\begin{tabular}{llrrrr}
\toprule
model & vs & mean\_diff & tstat & pvalue & n \\
\midrule
lasso & stockformer & 0.0319 & 1.8616 & 0.0639 & 236 \\
ridge & stockformer & 0.0187 & 1.4215 & 0.1565 & 236 \\
lightgbm & stockformer & 0.0178 & 1.6267 & 0.1052 & 236 \\
xgboost & stockformer & 0.0214 & 1.7157 & 0.0875 & 236 \\
momentum & stockformer & -0.0093 & -0.5822 & 0.5610 & 236 \\
\bottomrule
\end{tabular}
\end{table}


La mayor diferencia corresponde al Lasso, con $+0{,}0319$ de IC frente a Stockformer,
$t=1{,}86$ y $p=0{,}064$. Este $p$-valor no llega al umbral convencional de $0{,}05$, de
modo que no se puede rechazar la hipótesis nula de igualdad de IC al nivel del 5\,\%: no
cabe afirmar que el Lasso bata a Stockformer de forma estadísticamente significativa. Lo
que sí se observa es que la dirección del efecto es la misma en toda la escalera. XGBoost
supera a Stockformer en $+0{,}0214$ ($t=1{,}72$, $p=0{,}088$), LightGBM en $+0{,}0178$
($t=1{,}63$, $p=0{,}105$) y Ridge en $+0{,}0187$ ($t=1{,}42$, $p=0{,}157$). Todos los
modelos simples superan a Stockformer en IC con el mismo signo, con $p$-valores entre
$0{,}06$ y $0{,}16$; ninguno cruza el umbral. Los contrastes más conservadores de la
Tabla~\ref{tab:power-signtest} (prueba de signos y Wilcoxon) lo confirman, con intervalos
de confianza que incluyen el cero en todos los casos: la superioridad es direccional, no
significativa. El único control que no supera a Stockformer es el momentum ($-0{,}0093$,
$p=0{,}561$), que por construcción no es un modelo sino ruido.

\begin{figure}[htbp]
  \centering
  \includegraphics[width=0.85\linewidth]{ic_with_ci}
  \caption{IC medio por modelo con su intervalo de confianza al 95\,\% construido por
  remuestreo de días de negociación (2.000 réplicas), sobre la ventana de prueba de la
  escalera de complejidad. La definición del IC y el procedimiento de remuestreo se
  describen en la sección~\ref{sec:met-harness}.}
  \label{fig:ic-with-ci}
\end{figure}

La Figura~\ref{fig:ic-with-ci} acompaña esta lectura. El intervalo de Stockformer solapa
el cero y queda en el extremo inferior de la escalera. Los modelos simples están en
territorio positivo, pero sus intervalos también rozan el cero, lo que es coherente con un
mercado eficiente, donde de precio y volumen públicos no se extrae un IC transversal
grande y estable. El punto no es que los modelos simples tengan una señal potente, que no
la tienen, sino que el modelo de un millón de parámetros no mejora a los de unos pocos y
es el único con IC negativo.

Esa consistencia direccional no debe sobreinterpretarse. Que los cinco modelos superen a
Stockformer con el mismo signo llama la atención, pero no son cinco evidencias
independientes: los tres lineales son casi colineales y los dos modelos de árboles
comparten familia, de modo que el número efectivo de enfoques independientes es del orden
de dos. El contraste menos afectado por esa dependencia es la prueba de signos temporal,
que mide para cada modelo en qué fracción de los días su IC supera al de Stockformer. La
Tabla~\ref{tab:power-signtest} la recoge, junto con el contraste de Wilcoxon, el intervalo
de confianza de la diferencia media de IC y el efecto mínimo detectable. El mejor modelo,
el Lasso, supera a Stockformer en 130 de 236 días ($55{,}1\%$), con un $p$-valor de la
prueba de signos de $0{,}134$ a dos colas y un Wilcoxon de $0{,}108$; el resto queda entre
el $50{,}8\%$ y el $53{,}8\%$. Ningún modelo cruza el umbral del 5\,\% por la unidad de
observación correcta.

\begin{table}[htbp]
\centering
\caption{Contrastes direccionales del IC diario frente a Stockformer (comparaci\'on
pareada por d\'ia, $n=236$ d\'ias comunes): porcentaje de d\'ias con IC superior al de
Stockformer, $p$-valor de la prueba de signos (binomial exacta, dos colas), $p$-valor de
Wilcoxon, intervalo de confianza al 95\,\% de la diferencia media de IC ($\Delta$IC) y
efecto m\'inimo detectable (MDE) a una potencia del 80\,\% y nivel del 5\,\%.}
\label{tab:power-signtest}
\small
\begin{tabular}{lrrrrr}
\toprule
Modelo & \% d\'ias $>$SF & $p$ signos & $p$ Wilcoxon & $\Delta$IC [IC 95\%] & MDE$_{80}$ \\
\midrule
Ridge & 50.8 & 0.845 & 0.254 & $+0.0187$ [-0.0072, +0.0446] & 0.0369 \\
Lasso & 55.1 & 0.134 & 0.108 & $+0.0319$ [-0.0019, +0.0657] & 0.0481 \\
ElasticNet & 53.8 & 0.268 & 0.183 & $+0.0251$ [-0.0073, +0.0575] & 0.0461 \\
LightGBM & 53.8 & 0.268 & 0.164 & $+0.0178$ [-0.0038, +0.0393] & 0.0306 \\
XGBoost & 52.5 & 0.474 & 0.200 & $+0.0214$ [-0.0032, +0.0459] & 0.0349 \\
\bottomrule
\end{tabular}
\end{table}


Esta falta de significancia se entiende mejor a través de la potencia estadística. Con
$n=236$ días y la desviación típica observada de las diferencias de IC, el efecto mínimo
detectable a una potencia del 80\,\% y un nivel del 5\,\% está, según el modelo, entre
$0{,}031$ y $0{,}048$ en escala de diferencia de IC. Todas las diferencias observadas
(entre $0{,}018$ y $0{,}032$) caen por debajo de su propio efecto mínimo detectable, y el
intervalo de confianza al 95\,\% de cada una incluye el cero. El estudio está, por diseño,
infradimensionado para resolver efectos del tamaño que este régimen permite; la falta de
significancia individual refleja una señal débil medida sobre una muestra corta y
autocorrelacionada, no necesariamente la ausencia de todo efecto. Esta lectura sigue el
estándar de \textcite{harvey2016crosssection}, que pide estadísticos $t$ del orden de 3
para declarar un factor genuino bajo multiplicidad de contrastes; con valores $t$ en torno
a $1{,}5$--$1{,}9$, lo que hay es evidencia direccional, no un hallazgo demostrado.

La magnitud observada es, además, la esperable, y eso funciona como control de
verosimilitud. La banda de ruido que fijan los controles ($\pm 0{,}013$) sirve de
referencia, ya que el Lasso ($+0{,}0238$) la supera y Stockformer ($-0{,}003$) no. Un IC
sostenido por encima de $0{,}05$ en un universo tan arbitrado como el S\&P 500 sería un
indicio de fuga de datos. La auditoría de la sección~\ref{sec:met-leakage} confirma que no
hay fuga que perjudique a Stockformer; la única detectada juega a su favor y aun así su IC
es negativo. El hallazgo no se apoya, por tanto, en un único $p$-valor ni en la unanimidad
entre modelos correlacionados, sino en tres lecturas que apuntan a lo mismo: el IC de
Stockformer indistinguible de cero, la dirección consistente aunque no concluyente de los
modelos simples, y una potencia que sitúa todos los efectos dentro de lo que el régimen
permite resolver.

\section{Modos de fallo de la transferencia}\label{sec:p1-fallo}

Los resultados anteriores dejan abierta la pregunta de por qué no aparece la señal. Para
responderla hay que pasar de la constatación empírica a asociar cada observación a un
mecanismo documentado en la literatura. Hay cuatro modos de fallo que, actuando juntos,
explican la pérdida de señal al cruzar del mercado chino al S\&P 500.

\paragraph{(a) Eficiencia de mercado.} Stockformer fue calibrado sobre los índices
CSI~300/500, en un mercado dominado por inversores minoristas, con ineficiencias
persistentes y patrones de comportamiento que un modelo no lineal puede explotar
\autocite{liu2019size}. El S\&P 500 es uno de los mercados más eficientes y arbitrados del
mundo, dominado por inversores institucionales. Una señal que captura ineficiencias
minoristas en China encuentra poco que explotar en el S\&P 500, donde el arbitraje la ha
eliminado. Este es el modo de fallo más básico, y es coherente con que ningún modelo de la
escalera, ni simple ni profundo, supere un IC de $0{,}025$.

\paragraph{(b) Insuficiencia de la señal de precio y volumen.} Las variables de
Stockformer son esencialmente de precio y volumen (Alpha360 es ese tipo de información con
retardos). En un mercado eficiente, la información del historial de precio y volumen está,
en buena medida, ya incorporada al precio. La propia escalera lo confirma, ya que todos los
modelos ven el mismo panel de 375 variables, y los árboles, pese a su mayor capacidad, no
superan a los lineales sobre esos mismos insumos. El problema no está en la cantidad de
variables, sino en que esa clase de información no basta a este horizonte en este mercado.

\paragraph{(c) Microestructura.} El mercado chino impone límites diarios de fluctuación de
precio. Estos límites introducen una autocorrelación mecánica en los retornos (una acción
que toca el límite tiende a continuar al día siguiente) que infla el IC alcanzable por un
modelo de precio y volumen. El S\&P 500 no tiene estos límites, de modo que parte del IC
que Stockformer reporta en su mercado de origen no procede de poder predictivo genuino,
sino de un artefacto de microestructura que no existe en el mercado de destino.

\paragraph{(d) Dimensionalidad efectiva.} En cada corte temporal, el panel ofrece del
orden de 400 observaciones transversales efectivas frente a los del orden de un millón de
parámetros del modelo. Ajustar un modelo de ese tamaño sobre una sección transversal tan
estrecha tiende al sobreajuste, porque el modelo tiene capacidad para memorizar ruido
idiosincrático en vez de aprender señal generalizable. Esto explica por qué el modelo más
complejo es el peor (Sección~\ref{sec:p1-ladder}) y por qué un Lasso de cinco parámetros,
con su fuerte sesgo regularizador, generaliza mejor.

\paragraph{Contraste con los resultados reportados en la literatura.} La literatura
reciente reporta IC bastante más altos para arquitecturas profundas, pero siempre sobre
mercados distintos del S\&P 500. \textcite{fan2024stockmixer} reportan para StockMixer un
IC del orden de $0{,}041$ sobre NASDAQ, y \textcite{li2024master} un IC del orden de
$0{,}064$ para MASTER sobre CSI~300. Estas cifras, muy por encima de cualquier valor de la
escalera de este trabajo, no contradicen el hallazgo, pues son resultados obtenidos en sus
mercados de evaluación, que no tienen por qué trasladarse al S\&P 500, tal como sugieren
los cuatro modos de fallo anteriores. La transferibilidad entre mercados de los modelos de
aprendizaje profundo financiero no puede darse por hecha a partir de su desempeño en el
mercado de calibración; debe demostrarse empíricamente, y en el caso de esta
implementación de Stockformer sobre el S\&P 500 la evidencia disponible no la muestra.

En conjunto, los cuatro modos de fallo no son alternativas excluyentes, sino mecanismos
que se suman: un mercado más eficiente (a) reduce la señal explotable; la señal que queda
es de una clase de información insuficiente a este horizonte (b); parte del IC de origen
era un artefacto de microestructura inexistente en el destino (c); y la escasez de
observaciones transversales frente a la capacidad del modelo (d) convierte cualquier señal
débil en sobreajuste. El resultado neto es un IC de $-0{,}003$. Este diagnóstico cierra la
Parte~I y abre la pregunta del Capítulo~\ref{cap:parte2}. Si las redes complejas no
aportan señal en este régimen, de dónde sale el retorno neto que sí se observa.
